\chapter{算法实现与可证伪实验}
\label{chap:implementation}

\section{把外层目标与内层优化器分开实现}
实现采用四个独立对象：行为版本 $b_k$、镜像 anchor $\bar\pi_k$、当前 actor $p_\theta$、固定教师 $T$。
在 hard sync 且 fresh rollout 的基准中，前两个对象数值相同；在 replay、delayed sync 或 EMA 中，它们一般不同。
每条 response 保存生成版本、prompt、prefix、行为 log-prob 和采样配置。
anchor log-prob 必须用本轮冻结的 anchor 重算；行为 log-prob 只有在明确的测度校正项中使用。
温度、top-$p$、截断和 EOS 处理改变真实采样分布，保存模型 softmax log-prob 不能替代保存截断后分布。

每个 outer step 的顺序为：先同步本轮 anchor，随后采样或组合 replay，计算并停止梯度的教师信号，最后在该固定子问题上做 inner 更新。
一旦同步 anchor，旧残差的零点也随之移动，必须重新计算 $\log T-\log\bar\pi_k$。
只更新参数副本而继续复用旧 reward，得到的是另一条动力学。
EMA 则必须指定在 actor 更新前还是更新后执行，以及周期按 minibatch、rollout 还是 outer step 计数。
这些选择直接进入第~\ref{chap:ema}~章的矩阵，无法用同一个“同步率”概括。

优化器状态也需要固定语义。保留 Adam 动量会使下一轮的首个 inner 更新依赖上一轮目标。
清空动量给出更接近独立近端求解的基准，保留动量给出带记忆的求解器；两种实现应各自命名并记录。
理论中 $q^M$ 的收缩率对应固定二次模型的梯度下降，不能直接赋给任意 Adam 设置。

\section{Fisher 距离控制与数值求解}
令 $g$ 为当前镜像 surrogate 的参数梯度，$F$ 为同一 prefix 测度下的 Fisher 矩阵。
求解阻尼系统
\[
 (F+\lambda I)d=g,\qquad \lambda>0.
\]
共轭梯度无需显式存储矩阵，可用 Fisher-vector product 计算每次迭代。
当目标是最大化 surrogate，候选更新写成
\[
 \Delta\theta=\sqrt{\frac{2\epsilon_{\mathrm{NG}}}{d^\top Fd}}\,d.
\]
若 $d^\top Fd=0$，停止该方向更新；该分母为零表明所估 Fisher 无法识别此方向，继续放大步长没有几何依据。
阻尼系统与距离归一化使用的矩阵不同，因此不能将分母随意替换为 $g^\top d$。
无阻尼且精确求解时，两者相等；有阻尼时
\[
 g^\top d=d^\top Fd+\lambda\lVert d\rVert^2.
\]

以冻结 batch 的实际 $\KL(p_{\theta+\Delta\theta}\Vert p_\theta)$ 验证二阶预测，超预算时回溯缩短步长。
这控制的是 inner 求解路径；$\eta$ 控制镜像目标的位置。
固定 $\eta$、改变 $\epsilon_{\mathrm{NG}}$，理论零点不变；固定几何半径、改变 $\eta$，零点改变。
有限训练预算下，两者都会改变最终 student，因此应同时报告目标位置和追赶误差。

\section{与理论匹配的停止条件}
对于 tabular 全支持基准，可以直接计算
\[
 E_k=\sup_s\osc_a\log\frac{p_{k,M}(a\mid s)}{G_{\eta_k}(\bar\pi_k)(a\mid s)}.
\]
按 $E_k\le c\eta_k$ 停止，控制固定误差地板；要得到逐渐精确的 teacher 收敛，则要求 $E_k/\eta_k\to0$。
神经网络中只能在验证 prefix 集合估计该量，不能将有限集合的最大值当成所有 prefix 的最大值。
同时报告训练 residual、验证 target KL、深度分层覆盖率，识别过拟合和缺少访问的状态。

平方残差均值、梯度范数和 projective 误差具有不同量纲。
若残差平方的阈值随 $\eta_k^2$ 缩放，还需要 support 下界和范数转换常数，才能推出 $E_k=O(\eta_k)$。
简单地让 loss 小于一个固定数字，在小 $\eta$ 时会让求解误差支配镜像信号。

\section{四项研究假设及其反证}
\subsection{残差回归与序列 IS 的方差}
在相同 replay、相同 student、相同目标和相同 token 预算下，比较序列 IS 梯度与最优性残差梯度的二阶矩。
第~\ref{chap:opd}~章的条件期望结论只保证可用 Rao--Blackwell 化删除无关的 suffix 权重；
残差回归改变优化目标，方差更小必须通过估计器分析和实验共同检验。
按 response 长度报告梯度 trace、最大样本贡献与有效样本量
\[
 \mathrm{ESS}=\frac{(\sum_i w_i)^2}{\sum_i w_i^2}.
\]
较小方差同时伴随较大偏差时，应报告两者和最终目标误差，不能只展示 ESS。

\subsection{局部目标与序列目标的长期影响}
在 dense teacher log-ratio 和 sparse verifier reward 两个 regime 分别比较。
逐 prefix 的几何 teacher 插值有局部闭式解，完整序列 Gibbs 还包含 continuation partition function。
因此在 sparse reward 下，局部化后的 value 或 advantage 误差必须单独测量。
对相同 prompt 测试浅层和深层 prefix，并记录 suffix 成功率；若局部 KL 下降而深层成功率恶化，说明局部目标与长程 credit 的差异主导结果。

\subsection{镜像步长与实际几何步长的可分辨性}
先固定 $\eta$ 扫描 $\epsilon_{\mathrm{NG}}$，再固定半径扫描 $\eta$。
记录目标分布、实际 inner KL、target KL 与更新次数。
当足够多 inner 更新后不同半径仍给出不同零点，应检查投影容量、测度变化和停止梯度实现。
预测 KL 与实际 KL 长期不符，则 Fisher 估计或局部二阶近似不适合作为控制器。

\subsection{旧数据的覆盖与测度误差}
使用 fresh-only、replay-only 和 mixed-prefix 三组数据，保持教师和 anchor log-prob 重算一致。
先在完整 tabular support 中隔离纯测度效应，再在有限神经网络中加入未访问 prefix。
第~\ref{chap:guarantees}~章给出有界 integrand 的测度替换误差；零残差共享解并不意味着有限参数投影与测度无关。
当 coverage 相同而 replay 性能仍明显下降，应测量加权投影误差、行为概率记录和版本漂移，不能预先归因为 coverage。

\section{实验矩阵与公平预算}
\begin{table}[htbp]
\centering
\begin{tabular}{p{0.14\textwidth}p{0.79\textwidth}}
\toprule
实验轴 & 设置与对应判据\\\midrule
目标 & SFT、vanilla reverse-KL OPD、sequence PMD-part、PMD-mean、BPO、prefix MirrorOPD；每项明确优化目标和 prefix 测度。\\
数据 & fresh rollout、1/4/16/64 版本滞后、不同 replay 混合比例；同时报告 wall time 与有效 token。\\
测度修正 & 无 IS、action IS、per-decision IS、sequence IS、残差回归；区分无偏修正与换目标。\\
inner 求解 & Adam、KL 回溯、Fisher/NG；比较固定更新数、固定 token 数与同 residual 三种预算。\\
镜像强度 & $\eta\in\{0.1,0.3,0.5,0.7,1\}$；另保留原温度网格 $\tau\in\{0.25,0.5,1,2,4\}$，其中 $\eta=\tau^{-1}>1$ 标注为外推区间。\\
同步 & hard、delayed、参数 EMA、logit EMA；记录实际同步周期、EMA 系数及 pre/post-update 顺序。\\
信号 & dense teacher、终局 verifier、advantage regression；教师质量与支持下界作为独立轴。\\
诊断 & reward、pass@k、双向 KL、Fisher 距离、ESS、$\E[(\log W)^2]$、梯度 trace、prefix coverage、inner residual 和 actor--anchor lag。\\
\bottomrule
\end{tabular}
\end{table}

先在有限动作、有限前缀树上实现精确求解和有限 inner 更新。
精确 hard sync 的振幅曲线应满足 $D_k=|1-\eta|^kD_0$，$\eta=1$ 在一次全支持投影后为零。
有固定 inner 误差时应出现 $O(\varepsilon/\eta)$ 地板；EMA 模态应吻合矩阵谱半径及稳定边界。
这一步检验数学模型与代码语义一致，再进入神经网络容量和采样误差实验。

最终比较使用相同总 teacher 查询数、生成 token 数和训练 token 数，并独立报告 wall time。
较小 $\eta$ 所需 outer 次数增加，不能用“每轮更稳”直接推出总训练更快。
EMA 需要用发散率、可用 actor 步长和达到同一 target KL 的成本证明收益；仅展示较平滑 loss 曲线不足以说明收益。

\section{有限时域同步审计}
为了把上述递推和实现语义分开检验，目录
\texttt{MirrorOPD/experiment/outer\_dynamics\_qa/} 提供了一个可复现的有限时域
二动作实验。环境有 $H=6$ 个独立 prefix 状态，每个状态有两个动作；每个 outer
冻结 $\pi_k$，从它采样，执行固定数量的 inner 更新，然后 hard-sync 得到
$\pi_{k+1}$。教师由有限时域 reward 的 Boltzmann 分布构造，因而可以直接计算
精确的 $G_{\eta_k}(\pi_k)$、实际策略与 target 的 KL，以及每轮 reward gain。

实验同时运行 PMD-Part、PMD-Mean、BPO 和 MDPO 四个 inner 变体，并在相同随机种子、
样本和学习率下比较固定 $\eta$ 与从 $0.2$ 线性增加到 $0.7$ 的调度。运行命令为：
\begin{verbatim}
uv run --with numpy --with matplotlib python \\
  科研/pmd-theory/MirrorOPD/experiment/outer_dynamics_qa/run_finite_horizon.py
\end{verbatim}
输出 CSV 保存每个 outer 的 target displacement 和 finite-inner residual；审计文件
记录了 seed、预算和最终数值。图~\ref{fig:finite-outer-audit} 展示 reward、
策略到当轮 target 的 KL，以及单轮 reward gain。

实现层面的 PMD-Mean 使用 rollout batch 的中心化 reward 均值作为估计量，因此它是
定理~\ref{thm:four-inner-pmdmean-lambert} 中总体 PMD-Mean 目标的有限样本近似；代码
没有在每个历史上显式求解 Lambert-$W$ 归一化方程。实验中的 PMD-Mean 曲线应同时包含
均值估计误差和该总体 target bias，不能把它解释为 exact mirror target 的纯优化误差。

\begin{figure}[htbp]
\centering
\includegraphics[width=0.98\textwidth]{../MirrorOPD/experiment/outer_dynamics_qa/finite_horizon_outer_dynamics.pdf}
\caption{有限时域 MirrorOPD 同步审计。四种 inner 变体共享冻结 anchor、teacher、
采样批次和优化预算；线性调度只在 outer 边界改变 $\eta_k$。}
\label{fig:finite-outer-audit}
\end{figure}

该实验的作用是验证动力学，而不是替代神经网络实验。具体地，
\texttt{target\_kl\_anchor} 检验精确镜像位移，\texttt{kl\_policy\_target}
检验有限 inner 误差，二者分别对应定理~\ref{thm:projective} 和
定理~\ref{thm:inexact} 中的两个量。若增加 inner steps 后前者保持不变而后者下降，
说明实现确实把 $\eta$（目标位置）和 inner 学习率（求解精度）分开了。
